\providecommand{\bW}{\mathbb{W}}

\subsection{Artin--Schreier--Witt Theory}
\phantomsection\label{subsection:witt_vector_prerequisites}

Artin--Schreier theory, together with the ramification formula of
\Cref{theorem:artin_schreier_ramification}, gives complete control over the cyclic
extensions of degree $p$ of a local field of characteristic $p$. To treat the groups
$G = \Z/p^r\Z$ we need the analogous control over cyclic extensions of degree $p^r$,
which is provided by Witt's generalization: the field $K$ is replaced by the ring
$W_r(K)$ of $p$-typical Witt vectors of length $r$, the operator $\wp(u) = u^p - u$ by
its Witt-vector analogue $\wp = F - \mathrm{id}$, and the classical ramification-break
formula by the Schmid--Witt formula. The underlying Witt-ring algebra is collected in
\Cref{appendix:witt_vector_algebra}; this subsection retains the ASW classification,
the identities used in the proof, and the ramification statements. Our references for
the general theory are \cite[II, \S6]{serreLF} and \cite{Rabinoff2014Witt}.

Throughout, $p$ is a fixed prime, and $r \geq 1$ is an integer. All the Witt vectors
appearing in this thesis are $p$-typical and of finite length.


\paragraph{Working Witt-vector notation.}
For an \(\F_p\)-algebra \(A\), write \(W_r(A)=A^r\) for the ring of
\(p\)-typical Witt vectors of length \(r\), with Witt addition and subtraction
denoted by \(\oplus\) and \(\ominus\). We use the componentwise Frobenius \(F\),
the Verschiebung \(V\), and the truncation \(t\). The construction from Witt
polynomials, the ghost-map arguments, and the proofs of the identities among
these operators are given in \Cref{appendix:witt_vector_algebra}. In particular,
we use the triangular form of Witt addition
(\Cref{eq:witt_addition_shape}), the identities
\[
tW_r(A)=W_{r-1}(A),\qquad
\ker(t)=V^{r-1}A,\qquad
FV=VF=p,
\]
and the canonical identification \(W_r(\F_p)\cong\Z/p^r\Z\)
(\Cref{prop:witt_multiplication_by_p,example:witt_of_Fp}).

\begin{definition}\label{definition:asw_operator}
The \textbf{Artin--Schreier--Witt operator} on $W_r(A)$, $A$ an $\F_p$-algebra, is
\[
\wp \;:=\; F - \mathrm{id},
\qquad
\wp\,\vec u \;=\; F\vec u \ominus \vec u .
\]
For $r = 1$ this is the classical operator $\wp(u) = u^p - u$ of
\Cref{theorem:artin_schreier_ramification}.
\end{definition}

Since $F$ is a ring endomorphism, $\wp$ is an endomorphism of the additive group
$(W_r(A), \oplus)$:
$\wp(\vec u \oplus \vec v) = \wp\vec u \oplus \wp\vec v$. Moreover $\wp$ commutes with
$W_r(\sigma)$ for every ring homomorphism $\sigma$ (both $F$ and the ring operations
are natural), and with $V$ and $t$ by \Cref{prop:witt_multiplication_by_p}.

The following bookkeeping lemma is elementary but does the essential work of reducing
statements about $W_r$ to statements about $W_{r-1}$ and about $A$ itself: the
truncation $t$ and the inclusion $V^{r-1}$ of its kernel slice $W_r$ into lower-length
pieces, compatibly with $\wp$, and \emph{addition of an element of the kernel is
carry-free}.

\begin{lemma}[Witt bookkeeping]\label{lemma:witt_bookkeeping}
Let $A$ be an $\F_p$-algebra.
\begin{enumerate}[label=(\roman*)]
    \item $t \colon W_r(A) \to W_{r-1}(A)$ is a ring homomorphism commuting with $F$,
    hence with $\wp$; its kernel is $V^{r-1} A$.
    \item $V^{r-1} \colon A \to W_r(A)$ is additive, and
    $\wp \circ V^{r-1} = V^{r-1} \circ \wp$, where on the left $\wp$ is the operator
    on $W_r(A)$ and on the right the classical $\wp(c) = c^p - c$ on $A$.
    \item (carry-free top-level addition) For all $\vec a \in W_r(A)$ and $c \in A$,
    \[
    \vec a \oplus V^{r-1}(c) \;=\; (a_0, \dots, a_{r-2},\, a_{r-1} + c).
    \]
\end{enumerate}
\end{lemma}

\begin{proof}
(i) and (ii) restate parts of \Cref{prop:witt_multiplication_by_p} (for (ii): $V^{r-1}$
is additive and commutes with the componentwise $F$, hence with
$\wp = F - \mathrm{id}$).
(iii) By \Cref{lemma:ghost_principle} it suffices to verify the identity over
$\Z[a_0, \dots, a_{r-1}, c]$, where the ghost map is injective. Put
$\vec b := (0, \dots, 0, c)$ and $\vec s := \vec a \oplus \vec b$. For $n < r-1$ we
have $w_n(\vec b) = 0$, so $w_n(\vec s) = w_n(\vec a)$, and by induction on $n$
(solving for $s_n$, using that $p$ is a non-zero-divisor) $s_n = a_n$ for
$n \leq r-2$. For $n = r-1$, $w_{r-1}(\vec b) = p^{r-1} c$ gives
$p^{r-1} s_{r-1} = p^{r-1} a_{r-1} + p^{r-1} c$, i.e.\ $s_{r-1} = a_{r-1} + c$.
\end{proof}

\begin{rem*}
\Cref{lemma:witt_bookkeeping}(iii) fails for the addition of vectors supported in
lower positions: in general the carry polynomials $C_n$ of
\Cref{eq:witt_addition_shape} are non-zero, and Witt addition is \emph{not}
componentwise. In particular, operations performed componentwise on a Witt vector ---
such as discarding the regular parts of Laurent-series components --- do \emph{not}
in general respect Witt addition, and will be carefully avoided.
\end{rem*}

\subsubsection*{The Artin--Schreier--Witt isogeny and its torsors}

The componentwise ring operations make $\A^r_{\F_p}$ a ring scheme, the
\textbf{Witt vector group scheme} $\bW_r$: its underlying scheme is
$\A^r_{\F_p} = \Spec \F_p[T_0, \dots, T_{r-1}]$, with the group law given by the
addition polynomials $S_n$, so that $\bW_r(A) = W_r(A)$ functorially. The operator
$\wp = F - \mathrm{id}$ is an endomorphism of the group scheme $\bW_r$, and it is the
Witt-vector instance of a \emph{Lang isogeny}: by \Cref{example:witt_of_Fp} and the
lemma below, it sits in a short exact sequence of group schemes over $\F_p$ for the
\'etale topology
\begin{equation}\label{eq:asw_lang_sequence}
0 \longrightarrow \Z/p^r\Z \cong W_r(\F_p)
\longrightarrow \bW_r \xrightarrow{\;\wp\;} \bW_r \longrightarrow 0 .
\end{equation}

The corresponding covers admit the following completely explicit description, which
is the form in which we use them. For $\vec g \in W_r(A)$ we write
$A[\vec X]/(\wp \vec X \ominus \vec g)$, $\vec X = (X_0, \dots, X_{r-1})$, for the
quotient of $A[X_0, \dots, X_{r-1}]$ by the ideal generated by the $r$ components of
the Witt vector $\wp \vec X \ominus \vec g \in W_r(A[X_0, \dots, X_{r-1}])$.

\begin{lemma}[ASW covers are \'etale torsors]\label{lemma:asw_etale_torsor}
Let $A$ be an $\F_p$-algebra and $\vec g \in W_r(A)$. Let
\[
B \;:=\; A[\vec X]\big/\bigl(\wp \vec X \ominus \vec g\bigr).
\]
Then $B$ is finite free over $A$ with basis
$\bigl\{\prod_{n} X_n^{a_n} : 0 \le a_n \le p-1\bigr\}$, and \'etale over $A$. The
translation action of $W_r(\F_p) \cong \Z/p^r\Z$,
\[
m \colon \vec X \longmapsto \vec X \oplus \underline m ,
\]
makes $\Spec B$ a $\Z/p^r\Z$-torsor over $\Spec A$.
\end{lemma}

\begin{proof}
By \Cref{eq:witt_addition_shape} and the componentwise formula for $F$, the $n$-th
defining equation has the triangular form
\[
X_n^p - X_n - g_n + C_n'(X_0, \dots, X_{n-1};\, g_0, \dots, g_{n-1}) \;=\; 0
\]
for universal integral polynomials $C_n'$. Triangularity gives the free basis; the
Jacobian matrix of the system is lower triangular with diagonal entries
$p X_n^{p-1} - 1 = -1$, a unit, so $B$ is \'etale over $A$. Finally, $\Spec B$ is the
pullback along $\vec g \colon \Spec A \to \bW_r$ of the covering
$\wp \colon \bW_r \to \bW_r$; by the triangular form just established (over
$A = \F_p[T_0,\dots,T_{r-1}]$, $\vec g = (T_0,\dots,T_{r-1})$), $\wp$ is a surjective
\'etale homomorphism of group schemes, and its kernel is
$\{\vec u : F \vec u = \vec u\} = W_r(\F_p)$ since $F$ is componentwise. A surjective
\'etale group homomorphism is a torsor under its kernel, and torsors pull back to
torsors; hence $\Spec B$ is a $W_r(\F_p)$-torsor over $\Spec A$.
\end{proof}

\subsubsection*{Artin--Schreier--Witt theory for fields}

Over a field, the sequence \Cref{eq:asw_lang_sequence} computes all cyclic extensions
of $p$-power degree, exactly as classical Artin--Schreier theory does for degree $p$.

\begin{theorem}[Artin--Schreier--Witt theory; {\cite[\S3--4]{Thomas2005}}]
\label{theorem:asw_theory}
Let $M \supseteq \F_p$ be a field, $M^{s}$ a separable closure, and
$\Gamma = \Gal(M^{s}/M)$. Then $\wp$ is surjective on $W_r(M^{s})$ with kernel
$W_r(\F_p)$, and there is a canonical isomorphism of abelian groups
\[
\delta \colon\; W_r(M)\big/\wp\, W_r(M) \;\xrightarrow{\;\cong\;}\;
\Hom_{\mathrm{cont}}\bigl(\Gamma,\, \Z/p^r\Z\bigr) \;=\; H^1\bigl(M, \Z/p^r\Z\bigr),
\qquad
[\vec g] \longmapsto \bigl(\sigma \mapsto \sigma(\vec u) \ominus \vec u\bigr),
\]
where $\vec u \in W_r(M^{s})$ is any solution of $\wp \vec u = \vec g$. Under the
correspondence between characters and torsors
(\Cref{cor:abelian_equivilance_fundamental_torsors}), the class $[\vec g]$ corresponds
to the $\Z/p^r\Z$-torsor
\[
\Spec\, M[\vec X]\big/\bigl(\wp \vec X \ominus \vec g\bigr)
\]
of \Cref{lemma:asw_etale_torsor}. Moreover, for a class of order $p^s$ in
$W_r(M)/\wp W_r(M)$, the fixed field $M(\vec u) := M(u_0, \dots, u_{r-1})$ of
$\ker \delta[\vec g]$ is a cyclic extension of $M$ of degree $p^s$; the torsor is
connected if and only if $[\vec g]$ has order $p^r$, in which case
$\Spec M[\vec X]/(\wp \vec X \ominus \vec g) \cong \Spec M(\vec u)$ is (the spectrum
of) a cyclic field extension of degree $p^r$. Every cyclic extension of $M$ of degree
dividing $p^r$ arises this way.
\end{theorem}

\begin{proof}[Sketch of proof]
Surjectivity of $\wp$ on $W_r(M^{s})$ and the computation of its kernel follow from
the triangular equations in the proof of \Cref{lemma:asw_etale_torsor}: each component
of a preimage is obtained by solving a separable equation $X_n^p - X_n = (\dots)$, and
$F\vec u = \vec u$ holds if and only if $u_n^p = u_n$ for all $n$, i.e.\
$\vec u \in W_r(\F_p)$. Given $\vec g$ and a solution $\wp\vec u = \vec g$, for
$\sigma \in \Gamma$ we have
$\wp(\sigma \vec u \ominus \vec u) = \sigma \vec g \ominus \vec g = 0$ (using
\Cref{eq:witt_invariants_componentwise} and that $W_r(\sigma)$ commutes with $\wp$),
so $\sigma \vec u \ominus \vec u \in W_r(\F_p)$; one checks it is a continuous
homomorphism in $\sigma$ (the kernel is fixed pointwise by $\Gamma$), independent of
the choices modulo nothing and of the representative $\vec g$ modulo $\wp W_r(M)$.
Injectivity: if $\delta[\vec g] = 0$ then $\vec u$ is $\Gamma$-invariant, hence lies
in $W_r(M)$ by \Cref{eq:witt_invariants_componentwise}, so
$\vec g = \wp\vec u \in \wp W_r(M)$. Surjectivity of $\delta$ is the additive
Hilbert~90: $H^1(\Gamma, W_r(M^{s})) = 0$, by induction on $r$ along the exact
sequence $0 \to M^{s} \xrightarrow{V^{r-1}} W_r(M^{s}) \xrightarrow{t} W_{r-1}(M^{s})
\to 0$ of $\Gamma$-modules (\Cref{lemma:witt_bookkeeping}(i)) and the classical case
$H^1(\Gamma, M^{s}) = 0$. The remaining statements follow from
\Cref{theorem:equivalence_fundamental_covers} applied to the character
$\delta[\vec g]$, whose image has the same order as the class $[\vec g]$. For details
see \cite[\S3--4]{Thomas2005} or \cite{Rabinoff2014Witt}.
\end{proof}

The order of a class is partially visible in its $0$-th component:

\begin{cor}\label{cor:asw_class_order_and_first_component}
Let $M \supseteq \F_p$ be a field and $\vec g \in W_r(M)$ with $g_0 = 0$. Then
$p^{r-1}[\vec g] = 0$ in $W_r(M)/\wp W_r(M)$; equivalently, the cyclic extension
attached to $[\vec g]$ has degree at most $p^{r-1}$. Consequently, if $[\vec g]$ has
order $p^r$ (e.g.\ if $\vec g$ defines a connected torsor), then \emph{every}
representative of the class has non-zero $0$-th component.
\end{cor}

\begin{proof}
By \Cref{eq:witt_p_to_r_minus_1}, $p^{r-1} \vec g = (0, \dots, 0, g_0^{p^{r-1}}) = 0$
already in $W_r(M)$.
\end{proof}

\begin{rem*}\label{rem:asw_local_totally_ramified}
When $M = k((x))$ with $k$ algebraically closed, every finite separable extension
$L/M$ is totally ramified: the residue field extension is trivial (as $k$ is
algebraically closed) and $\Oo_L$ is a complete discrete valuation ring, so
$[L : M] = ef = e$. Thus a class of order $p^r$ in $W_r(k((x)))/\wp$ corresponds to a
cyclic, \emph{totally ramified} extension of degree $p^r$.
\end{rem*}



The isobaricity property of the Witt addition polynomials used in the later
degree estimate is recorded, with its proof, in
\Cref{lemma:witt_addition_isobaric}.

\subsubsection*{Ramification breaks: the Schmid--Witt formula}

Finally we state the generalization of the ramification formula of
\Cref{theorem:artin_schreier_ramification} to Artin--Schreier--Witt extensions. As in
the classical case, the formula reads off the largest upper-numbering ramification
break from the pole orders of a defining Witt vector --- provided the representative
is normalized so that no interference between the levels can occur. The relevant
normalization is the following.

\begin{definition}\label{definition:reduced_witt_vector}
Let $K = k((x))$. For $g \in K$ put $m(g) := \max\bigl(0, -v_x(g)\bigr)$ (the pole
order of $g$, or $0$ if $g$ is regular). A vector
$\vec g = (g_0, \dots, g_{r-1}) \in W_r(K)$ is \textbf{reduced} if for every $j$,
either $m(g_j) = 0$ or $p \nmid m(g_j)$.
\end{definition}

\begin{theorem}[Schmid--Witt break formula]\label{theorem:schmid_witt_break_formula}
Let $K = k((x))$ with $k$ perfect of characteristic $p$, and let $L/K$ be a cyclic
totally ramified extension of degree $p^r$ whose Artin--Schreier--Witt class
(\Cref{theorem:asw_theory}) is represented by a \emph{reduced} vector
$\vec g \in W_r(K)$, with $m_j := m(g_j)$. Then the largest upper-numbering
ramification break of $L/K$ equals
\[
u \;=\; \max_{0 \le j \le r-1} \; p^{\,r-1-j}\, m_j .
\]
\end{theorem}

This is due, for finite residue fields, to Kanesaka--Sekiguchi
\cite{KanesakaSekiguchi1979} and Thomas \cite{Thomas2005} (building on Schmid
\cite{Schmid1937} for $r = 1$), and for arbitrary perfect residue fields to
Elder--Keating \cite{ElderKeating2025}. We use it as stated, in one direction only:
\emph{evaluating the formula on one reduced representative computes $u$}; no
minimization over representatives is needed.

\begin{rem*}
Reducedness is exactly what makes the formula exact. If two levels $j < j'$ with
$m_j, m_{j'} \neq 0$ satisfied $p^{r-1-j} m_j = p^{r-1-j'} m_{j'}$, then
$m_{j'} = p^{\,j'-j} m_j$ would be divisible by $p$, contradicting
\Cref{definition:reduced_witt_vector}; hence the maximum is attained at a unique
level, and no cancellation of breaks between levels can occur.
\end{rem*}

\begin{rem*}
For $r = 1$, a reduced vector is a single element $g \in K$ whose pole order $m_0$ is
zero or prime to $p$, and \Cref{theorem:schmid_witt_break_formula} recovers
\Cref{theorem:artin_schreier_ramification}: the largest (indeed the only) upper break
is $m_0$. The existence of reduced representatives --- and, over $K = k((x))$, of
reduced representatives whose components are polynomials in $x^{-1}$ without constant
term --- is not part of the general theory quoted here; it will be proved where it is
used, by induction on $r$ through the truncation maps of
\Cref{lemma:witt_bookkeeping}.
\end{rem*}
